% Pursuing Stacks (cotes 134-2 to 134-8): archive cote 134-7, site page 58 (PDF page 59), typescript page 500.
% typescript. Draft compared: ps6.tex:2858-2921.
% AI-transcribed by Claude Opus 5.5 (claude-opus-5-5, adaptive thinking, effort high), Message Batches API, batch ps4 (msgbatch_012Aoc1NuskATPsyBeKSYK7J), 2026-11-12; not yet checked by eye.
\scanpage{134-7/58}
\mspage{500}

hemispherical set $X_{\underline{*}}(k)$, handled as if it came from an actual $\infty$-groupoid
(\inserted{by} taking it's \xout{xxxx}\inserted[typed]{source} and target operations etc) even if it does not. (I con-%
fess I did not check that this process does correctly describe the homoto\unsure{-}%
py groups of a hemispherical set\hand{,} even \over{\unsure{with}}{with}out assuming it comes from
a\hand{n} $\infty$-groupoid or an $\infty$-stack, but never mind for the time being ...)
Among the interesting things still ahead (once we get a little \xout{xxxx} accus-%
tomed working with hemispherical complexes) is to try and understand how
these two descriptions relate to each other, which may be \deleted{a way}\inserted{one means} for \deleted{coming}
\deleted{to} a better understanding of the basic formula (9), in the context of
hemispheric schematic mod\over{\unsure{les}}{els}.


26.9.

119. After the last notes (of September 10) I was a little sick for a
few days, then I was taken by current tasks from professional and family
life, which left little leisure for mathematical reflection, except once
or twice for a couple of hours, by way of recreation. It would seem now
that in the days and weeks ahead, there will be more time to go on with
the notes, and I feel eager indeed to push ahead. Also I more or less
promised the publisher, Pierre Bérès, that a first volume would be ready
for the printer by the end of this calendar year, and I would like to
keep promise.

I still have to tie in with the reflections and happenings of the end
of last month, as I started upon with the last notes (of Sept. 10). Next
thing then to report upon is the ``coup de th\'e\hand{\^{a}}tre'' occuring through \xout{my}
the phone call to Luc Illusie. \over{\unsure{Then}}{When} I told him about what by then still
looked to me as the key assumption for a theory of schematization of homo-%
topy types, namely that the homology of $K(\hand{\pi},n)$ should be computable \hand{in}
terms of derived functors of the ``divided power algebra''-functor $\Gamma_{\underline{Z}}$ , he
at once felt rather sceptical, and later he called me back to tell me it
was definitely false. He could not give me an explicit counterexample for
$H_{i}(n,\underline{Z};\underline{Z})$ say with given n and i, rather he said that when suitably ``sta-%
bilizing'' the assumption I had in mind, it went against result\hand{s} of Larry
Breen on $\mathrm{Ext}^{i}$ functors of $\underline{G}_{a}$ with itself over prime fields $\underline{F}_{p}$ . I do'nt
know if I am going some day to dive into Illusie's argument and into Larry
Breen's results - \deleted{at any rate}\inserted{however}, even before I got Illusie's confirmation
that definitely my assumption was wrong, I convinced myself that at any
rate it was false for $n=1$ \over{,}{\unsure{.}} \inserted{This is} a non simply connected case and hence not

\margin{58}
